% Physical p.11 / Roman IX. Source line-break hyphens are retained here.
\IllusieSetPhysicalPageSize{450.6bp}{704bp}
\begin{illusiefrontpage}{illusie:I:intro:p011}{Introduction}
\phantomsection\hypertarget{illusie:I:intro}{}\hypertarget{illusie:I:intro:p011:p1}{}%
Le présent travail est consacré à la théorie du complexe cotangent sur les topos
annelés (globalisation du formalisme d'André ([1], [2]) et Quillen ([24], [25])),
et à des applications de celle-ci à quelques problèmes de déformation infi-\linebreak
nitésimale globale. Les plus intéressantes concernent des calculs d'obstructions à
la déformation de schémas en groupes plats commutatifs, calculs qui jouent un
rôle-clef dans les travaux récents de Grothendieck\textsuperscript{(1)} et
Messing\textsuperscript{(2)} sur les groupes de Barsotti-Tate.

\vspace{4bp}
\phantomsection\hypertarget{illusie:I:intro:p011:p2}{}%
Voici, brièvement, le contenu des divers chapitres.

\vspace{4bp}
\phantomsection\hypertarget{illusie:I:intro:p011:p3}{}%
Le chapitre I rassemble le matériel d'algèbre homologique nécessaire à la
construction et l'étude des propriétés du complexe cotangent. Nous introduisons la
notion de faisceau d'homotopie d'un faisceau simplicial et prouvons une variante
globale d'un théorème de Whitehead [9], à savoir que si un morphisme
$f:X\longrightarrow Y$ de faisceaux simpliciaux induit des isomorphismes sur les
faisceaux d'homotopie en tout degré, il en est de même du morphisme
$\mathbb Z(f):\mathbb Z(X)\longrightarrow\mathbb Z(Y)$ induit sur les
$\mathbb Z$-Modules simpliciaux libres engendrés. En ``inversant'' les flèches qui
indui-\linebreak sent des isomorphismes sur les faisceaux d'homotopie, on définit
diverses catégories ``dérivées'', généralisations communes des catégories dérivées
de Verdier [27] et des catégories homotopiques de Quillen [23]. En raison de son
importance pour la théorie du complexe cotangent, nous étudions plus spécialement
la catégorie dérivée des Modules sur un Anneau simplicial : triangles distingués,
$\operatorname{Ext}^{i}$, produit tensoriel dérivé. Le dernier paragraphe contient
la définition des foncteurs dérivés de certains foncteurs non additifs sur la
catégorie dérivée $D^{-}$ d'un topos annelé, et les analogues globaux des
isomorphismes de décalage de Quillen ([24] 7.21).

\IllusieNoteRule
\phantomsection\hypertarget{illusie:I:intro:p011:n1}{}%
\IllusieNote{1}{A. Grothendieck, \emph{Cristaux et modules de Dieudonné}, exposé au Congrès Internatio-\linebreak nal des Mathématiciens, Nice (1970), et Cours au Collège de France, 70-71 et 71-72.}
\phantomsection\hypertarget{illusie:I:intro:p011:n2}{}%
\IllusieNote{2}{W. Messing, \emph{The Crystals associated to Barsotti-Tate Groups, with Applications to Abelian Schemes}, Thesis, Princeton (1971), à paraître.}
\end{illusiefrontpage}
